\section{Anforderungen an die Gewichtungsfunktion}
\label{sec:anforderungen_an_die_gewichtungsfunktion}

Für die Konzeption der Gewichtungsfunktion ergeben sich aus dem Anwendungskontext der CTI-Reports folgende zentrale Anforderungen:

\begin{itemize}
    \item \textbf{Erzeugung fester Merkmalsvektoren:} Da die Anzahl der liefernden Knoten und damit der eingehenden CTI-Reports im operativen Betrieb kontinuierlich variiert, muss die Gewichtungsfunktion eine beliebige Anzahl von Eingangsmengen verarbeiten können. Ziel ist die Extraktion einer gleichbleibenden Informationsmenge (Fixed-Size Input) aus diesen variablen Eingangsdaten. Dies stellt sicher, dass der resultierende Merkmalsvektor unabhängig von der ursprünglichen Report-Menge eine einheitliche Struktur aufweist, was die notwendige Voraussetzung für die Performanz und Stabilität nachgelagerter ML-Algorithmen ist~\cite{zaheer2017deepsets, velickovic2018gat}.

    \item \textbf{Relevanzbasierte Gewichtung der Detektionen:} Um die Qualität der aggregierten Ergebnisse zu steigern, sollte das Modell die Relevanz der einzelnen CTI-Reports differenziert auf Basis der enthaltenen Einzeldetektionen bewerten. Diese Gewichtungsstrategie orientiert sich an modernen Attention-Mechanismen, wie sie in Kapitel~\ref{sec:vergleich_aggregationsmethoden} (Vergleich von Aggregationsmethoden) detailliert evaluiert wurden. Die Notwendigkeit einer solchen Attention-Mechanismusses ergibt sich aus der Anforderung, die Informationsdichte einzelner Quellen präzise abzubilden und signifikante CTI-Reports gegenüber weniger aussagekräftigen CTI-Reports zu priorisieren.

    \item \textbf{Asymptotisches Sättigungsverhalten der Gewichtung:} 
    Um eine übermäßige Dominanz einzelner CTI-Reports mit extrem hohen Detektionszahlen zu vermeiden, darf die Gewichtungsfunktion kein unbegrenztes Wachstum aufweisen. Gefordert wird ein fester Grenzwert (Asymptote), damit zusätzliche Detektionen ab einem gewissen Schwellenwert nur noch marginalen Einfluss auf das Gesamtgewicht haben. Dieser Mechanismus soll den Echo-Chamber-Effekt reduzieren (siehe Kapitel~\ref{sec:analyse_des_echo_chamber_effekts} (Analyse des Echo-Chamber-Effekts)) und verhindern, dass das Aggregationsergebnis durch redundante Meldungsmengen verzerrt wird.
\end{itemize}